
\documentclass{article}
\usepackage{colortbl}
\usepackage{makecell}
\usepackage{multirow}
\usepackage{supertabular}

\begin{document}

\newcounter{utterance}

\centering \large Interaction Transcript for game `clean\_up', experiment `2\_hard\_3obj\_hr', episode 3 with glm{-}5.2.
\vspace{24pt}

{ \footnotesize  \setcounter{utterance}{1}
\setlength{\tabcolsep}{0pt}
\begin{supertabular}{c@{$\;$}|p{.15\linewidth}@{}p{.15\linewidth}p{.15\linewidth}p{.15\linewidth}p{.15\linewidth}p{.15\linewidth}}
    \# & $\;$A & \multicolumn{4}{c}{Game Master} & $\;\:$B\\
    \hline

    \theutterance \stepcounter{utterance}  
    & & \multicolumn{4}{p{0.6\linewidth}}{
        \cellcolor[rgb]{0.9,0.9,0.9}{
            \makecell[{{p{\linewidth}}}]{
                \texttt{\tiny{[P1$\langle$GM]}}
                \texttt{Ja sam vaš voditelj igre, a vi igrate suradničku igru sa sljedećom mrežom kao pločom za igru:} \\
\\ 
\texttt{```} \\
\texttt{ 1234567} \\
\texttt{╔═╤═╤═══╗} \\
\texttt{║L│◌│◌◌◌║ 1} \\
\texttt{║C├─┤◌◌◌║ 2} \\
\texttt{╟─┤◌├─┬─╢ 3} \\
\texttt{║◌│◌│◌│◌║ 4} \\
\texttt{╟─┼─┘◌│◌║ 5} \\
\texttt{║◌│◌◌┌┴─╢ 6} \\
\texttt{║◌│◌◌│◌P║ 7} \\
\texttt{╚═╧══╧══╝} \\
\\ 
\texttt{```} \\
\\ 
\texttt{* Gornji rub prikazuje x{-}koordinate koje rastu udesno, a desni rub y{-}koordinate koje rastu prema dolje.} \\
\texttt{* Sljedeći objekti nasumično su postavljeni na vašoj mreži: 'C', 'L', 'P'.} \\
\\ 
\\ 
\texttt{Drugi igrač vidi varijantu ploče za igru, gdje su objekti postavljeni na različitim nasumičnim lokacijama. Ne možete vidjeti ploču drugog igrača, a ni on ne može vidjeti vašu.} \\
\\ 
\\ 
\texttt{**Cilj:**} \\
\\ 
\texttt{Oba igrača trebaju pomicati objekte na svojim pozadinama tako da identični objekti završe na istim koordinatama. Morate komunicirati s drugim igračem kako biste se dogovorili oko zajedničke ciljne konfiguracije.} \\
\\ 
\\ 
\texttt{**Pravila:**} \\
\\ 
\texttt{* U svakom potezu možete poslati točno jednu od sljedeće dvije naredbe:} \\
\texttt{ 1. `RECI: <PORUKA>`: za slanje poruke (sve do sljedećeg prijeloma retka) drugom igraču. Proslijedit ću je vašem partneru.} \\
\texttt{2. `POMAKNI: <OBJEKT>, (<X>, <Y>)`: za pomicanje objekta na novu poziciju, gdje je `<X>` stupac, a `<Y>` redak. Obavijestit ću vas je li vaš potez bio uspješan ili ne.} \\
\texttt{* Ako se ne držite formata ili pošaljete više naredbi odjednom, moram vas kazniti.} \\
\texttt{* Ako oba igrača skupe više od 8 kazni, oboje gubite igru.} \\
\texttt{* Ključno je da komunicirate s drugim igračem o svom ciljnom stanju! *Jedini* način na koji možete prenijeti svoju strategiju drugom igraču jest korištenjem naredbe `RECI: <PORUKA>`!} \\
\\ 
\\ 
\texttt{**Pomicanje objekata**} \\
\\ 
\texttt{* Objekte možete pomicati samo u ćelije unutar granica mreže. Ciljna ćelija mora biti prazna, tj. smije sadržavati samo simbol '◌'.} \\
\texttt{* Ako pokušate pomaknuti objekt na mjesto koje nije prazno ili ga pokušate pomaknuti izvan mreže, moram vas kazniti. Dobivate još jedan pokušaj.} \\
\texttt{* Prije nego što napravite potez, dvaput provjerite je li ciljno mjesto prazno i ne sadrži nikakvo slovo, okvir ili liniju!} \\
\\ 
\\ 
\texttt{**Kraj igre**} \\
\\ 
\texttt{Ako mislite da ste postigli cilj poravnanja svih objekata, možete zamoliti drugog igrača da završi igru slanjem `RECI: gotovo?`. Ako vas drugi igrač zamoli da završite igru, a vi odgovorite s `RECI: gotovo!`, igra će završiti.} \\
\\ 
\texttt{Oba igrača pobjeđuju ako igra završi unutar 12 rundi, pri čemu je jedna runda definirana kao to da oba igrača pošalju valjanu naredbu.} \\
\\ 
\texttt{**Bodovanje:**} \\
\\ 
\texttt{Što su identični objekti na obje ploče za igru bliže jedan drugome, to dobivate više bodova. Kazne smanjuju vaše bodove. Možete li srušiti rekord?} \\
\\ 
\\ 
\texttt{Molimo pošaljite poruku drugom igraču kako biste započeli igru!} \\
            }
        }
    }
    & \\ \\

    \theutterance \stepcounter{utterance}  
    & \multicolumn{4}{p{0.6\linewidth}}{
        \cellcolor[rgb]{0.9,0.9,0.9}{
            \makecell[{{p{\linewidth}}}]{
                \texttt{\tiny{[P1$\rangle$GM]}}
                \texttt{RECI: Pozdrav! Ja vidim svoje objekte na sljedećim koordinatama: L na (1,1), C na (1,2) i P na (7,7). Molim te, podijeli sa mnom koordinate tvojih objekata. Nakon toga možemo se dogovoriti o zajedničkom cilju. Predlažem da, ako je moguće, ciljamo pozicije koje su dosupne na obje ploče ili odaberemo neke jednostavne koordinate prema kojima ćemo oboje pomicati svoje objekte.} \\
            }
        }
    }
    & & \\ \\

    \theutterance \stepcounter{utterance}  
    & & & \multicolumn{2}{p{0.3\linewidth}}{
        \cellcolor[rgb]{0.9,0.9,0.9}{
            \makecell[{{p{\linewidth}}}]{
                \texttt{\tiny{[GM$|$GM]}}
                \texttt{RECI: Pozdrav! Ja vidim svoje objekte na sljedećim koordinatama: L na (1,1), C na (1,2) i P na (7,7). Molim te, podijeli sa mnom koordinate tvojih objekata. Nakon toga možemo se dogovoriti o zajedničkom cilju. Predlažem da, ako je moguće, ciljamo pozicije koje su dosupne na obje ploče ili odaberemo neke jednostavne koordinate prema kojima ćemo oboje pomicati svoje objekte.} \\
            }
        }
    }
    & & \\ \\

    \theutterance \stepcounter{utterance}  
    & & \multicolumn{4}{p{0.6\linewidth}}{
        \cellcolor[rgb]{0.9,0.9,0.9}{
            \makecell[{{p{\linewidth}}}]{
                \texttt{\tiny{[GM$\rangle$P2]}}
                \texttt{Ja sam vaš voditelj igre, a vi igrate suradničku igru sa sljedećom mrežom kao pločom za igru:} \\
\\ 
\texttt{```} \\
\texttt{ 1234567} \\
\texttt{╔═╤═╤═══╗} \\
\texttt{║◌│◌│◌◌◌║ 1} \\
\texttt{║◌├─┤◌L◌║ 2} \\
\texttt{╟─┤P├─┬─╢ 3} \\
\texttt{║◌│◌│◌│◌║ 4} \\
\texttt{╟─┼─┘C│◌║ 5} \\
\texttt{║◌│◌◌┌┴─╢ 6} \\
\texttt{║◌│◌◌│◌◌║ 7} \\
\texttt{╚═╧══╧══╝} \\
\\ 
\texttt{```} \\
\\ 
\texttt{* Gornji rub prikazuje x{-}koordinate koje rastu udesno, a desni rub y{-}koordinate koje rastu prema dolje.} \\
\texttt{* Sljedeći objekti nasumično su postavljeni na vašoj mreži: 'C', 'L', 'P'.} \\
\\ 
\\ 
\texttt{Drugi igrač vidi varijantu ploče za igru, gdje su objekti postavljeni na različitim nasumičnim lokacijama. Ne možete vidjeti ploču drugog igrača, a ni on ne može vidjeti vašu.} \\
\\ 
\\ 
\texttt{**Cilj:**} \\
\\ 
\texttt{Oba igrača trebaju pomicati objekte na svojim pozadinama tako da identični objekti završe na istim koordinatama. Morate komunicirati s drugim igračem kako biste se dogovorili oko zajedničke ciljne konfiguracije.} \\
\\ 
\\ 
\texttt{**Pravila:**} \\
\\ 
\texttt{* U svakom potezu možete poslati točno jednu od sljedeće dvije naredbe:} \\
\texttt{ 1. `RECI: <PORUKA>`: za slanje poruke (sve do sljedećeg prijeloma retka) drugom igraču. Proslijedit ću je vašem partneru.} \\
\texttt{2. `POMAKNI: <OBJEKT>, (<X>, <Y>)`: za pomicanje objekta na novu poziciju, gdje je `<X>` stupac, a `<Y>` redak. Obavijestit ću vas je li vaš potez bio uspješan ili ne.} \\
\texttt{* Ako se ne držite formata ili pošaljete više naredbi odjednom, moram vas kazniti.} \\
\texttt{* Ako oba igrača skupe više od 8 kazni, oboje gubite igru.} \\
\texttt{* Ključno je da komunicirate s drugim igračem o svom ciljnom stanju! *Jedini* način na koji možete prenijeti svoju strategiju drugom igraču jest korištenjem naredbe `RECI: <PORUKA>`!} \\
\\ 
\\ 
\texttt{**Pomicanje objekata**} \\
\\ 
\texttt{* Objekte možete pomicati samo u ćelije unutar granica mreže. Ciljna ćelija mora biti prazna, tj. smije sadržavati samo simbol '◌'.} \\
\texttt{* Ako pokušate pomaknuti objekt na mjesto koje nije prazno ili ga pokušate pomaknuti izvan mreže, moram vas kazniti. Dobivate još jedan pokušaj.} \\
\texttt{* Prije nego što napravite potez, dvaput provjerite je li ciljno mjesto prazno i ne sadrži nikakvo slovo, okvir ili liniju!} \\
\\ 
\\ 
\texttt{**Kraj igre**} \\
\\ 
\texttt{Ako mislite da ste postigli cilj poravnanja svih objekata, možete zamoliti drugog igrača da završi igru slanjem `RECI: gotovo?`. Ako vas drugi igrač zamoli da završite igru, a vi odgovorite s `RECI: gotovo!`, igra će završiti.} \\
\\ 
\texttt{Oba igrača pobjeđuju ako igra završi unutar 12 rundi, pri čemu je jedna runda definirana kao to da oba igrača pošalju valjanu naredbu.} \\
\\ 
\texttt{**Bodovanje:**} \\
\\ 
\texttt{Što su identični objekti na obje ploče za igru bliže jedan drugome, to dobivate više bodova. Kazne smanjuju vaše bodove. Možete li srušiti rekord?} \\
\\ 
\texttt{Drugi igrač je započeo igru slanjem ove poruke:} \\
\\ 
\texttt{"Pozdrav! Ja vidim svoje objekte na sljedećim koordinatama: L na (1,1), C na (1,2) i P na (7,7). Molim te, podijeli sa mnom koordinate tvojih objekata. Nakon toga možemo se dogovoriti o zajedničkom cilju. Predlažem da, ako je moguće, ciljamo pozicije koje su dosupne na obje ploče ili odaberemo neke jednostavne koordinate prema kojima ćemo oboje pomicati svoje objekte."} \\
\\ 
\texttt{Koja je vaša prva naredba?} \\
            }
        }
    }
    & \\ \\

    \theutterance \stepcounter{utterance}  
    & & & \multicolumn{4}{p{0.6\linewidth}}{
        \cellcolor[rgb]{0.9,0.9,0.9}{
            \makecell[{{p{\linewidth}}}]{
                \texttt{\tiny{[GM$\langle$P2]}}
                \texttt{RECI: Moji objekti su: L na (5,2), C na (4,5), P na (2,3). Predlažem cilj: L na (1,1), C na (1,2), P na (7,7). Možeš li tamo pomaknuti L i C?} \\
            }
        }
    }
     \\ \\

    \theutterance \stepcounter{utterance}  
    & & & \multicolumn{2}{p{0.3\linewidth}}{
        \cellcolor[rgb]{0.9,0.9,0.9}{
            \makecell[{{p{\linewidth}}}]{
                \texttt{\tiny{[GM$|$GM]}}
                \texttt{RECI: Moji objekti su: L na (5,2), C na (4,5), P na (2,3). Predlažem cilj: L na (1,1), C na (1,2), P na (7,7). Možeš li tamo pomaknuti L i C?} \\
            }
        }
    }
    & & \\ \\

    \theutterance \stepcounter{utterance}  
    & & \multicolumn{4}{p{0.6\linewidth}}{
        \cellcolor[rgb]{0.9,0.9,0.9}{
            \makecell[{{p{\linewidth}}}]{
                \texttt{\tiny{[P1$\langle$GM]}}
                \texttt{Vaša poruka je proslijeđena drugom igraču.} \\
\\ 
\\ 
\texttt{Trenutačno igrate rundu 1 od 12.} \\
\texttt{Zajedno ste prikupili 0 od 8 kazni.} \\
\texttt{Drugi igrač je poslao ovu poruku:} \\
\\ 
\\ 
\texttt{"Moji objekti su: L na (5,2), C na (4,5), P na (2,3). Predlažem cilj: L na (1,1), C na (1,2), P na (7,7). Možeš li tamo pomaknuti L i C?"} \\
\\ 
\\ 
\texttt{Koja je vaša sljedeća naredba?} \\
            }
        }
    }
    & \\ \\

    \theutterance \stepcounter{utterance}  
    & \multicolumn{4}{p{0.6\linewidth}}{
        \cellcolor[rgb]{0.9,0.9,0.9}{
            \makecell[{{p{\linewidth}}}]{
                \texttt{\tiny{[P1$\rangle$GM]}}
                \texttt{POMAKNI: L, (1, 1)} \\
            }
        }
    }
    & & \\ \\

    \theutterance \stepcounter{utterance}  
    & & & \multicolumn{2}{p{0.3\linewidth}}{
        \cellcolor[rgb]{0.9,0.9,0.9}{
            \makecell[{{p{\linewidth}}}]{
                \texttt{\tiny{[GM$|$GM]}}
                \texttt{POMAKNI: L, (1, 1)} \\
            }
        }
    }
    & & \\ \\

    \theutterance \stepcounter{utterance}  
    & & & \multicolumn{2}{p{0.3\linewidth}}{
        \cellcolor[rgb]{0.9,0.9,0.9}{
            \makecell[{{p{\linewidth}}}]{
                \texttt{\tiny{[GM$|$GM]}}
                \texttt{Kazna: (1,1) nije prazno, već sadrži 'L'.} \\
\\ 
\texttt{Zajedno ste prikupili 1 od 8 kazni.} \\
\texttt{Molim pokušajte ponovno!} \\
            }
        }
    }
    & & \\ \\

    \theutterance \stepcounter{utterance}  
    & & \multicolumn{4}{p{0.6\linewidth}}{
        \cellcolor[rgb]{0.9,0.9,0.9}{
            \makecell[{{p{\linewidth}}}]{
                \texttt{\tiny{[P1$\langle$GM]}}
                \texttt{Kazna: (1,1) nije prazno, već sadrži 'L'.} \\
\\ 
\texttt{Zajedno ste prikupili 1 od 8 kazni.} \\
\texttt{Molim pokušajte ponovno!} \\
            }
        }
    }
    & \\ \\

    \theutterance \stepcounter{utterance}  
    & \multicolumn{4}{p{0.6\linewidth}}{
        \cellcolor[rgb]{0.9,0.9,0.9}{
            \makecell[{{p{\linewidth}}}]{
                \texttt{\tiny{[P1$\rangle$GM]}}
                \texttt{POMAKNI: P, (7, 6)} \\
            }
        }
    }
    & & \\ \\

    \theutterance \stepcounter{utterance}  
    & & & \multicolumn{2}{p{0.3\linewidth}}{
        \cellcolor[rgb]{0.9,0.9,0.9}{
            \makecell[{{p{\linewidth}}}]{
                \texttt{\tiny{[GM$|$GM]}}
                \texttt{POMAKNI: P, (7, 6)} \\
            }
        }
    }
    & & \\ \\

    \theutterance \stepcounter{utterance}  
    & & & \multicolumn{2}{p{0.3\linewidth}}{
        \cellcolor[rgb]{0.9,0.9,0.9}{
            \makecell[{{p{\linewidth}}}]{
                \texttt{\tiny{[GM$|$GM]}}
                \texttt{Kazna: (7,6) nije prazno, već sadrži '─'.} \\
\\ 
\texttt{Zajedno ste prikupili 2 od 8 kazni.} \\
\texttt{Molim pokušajte ponovno!} \\
            }
        }
    }
    & & \\ \\

    \theutterance \stepcounter{utterance}  
    & & \multicolumn{4}{p{0.6\linewidth}}{
        \cellcolor[rgb]{0.9,0.9,0.9}{
            \makecell[{{p{\linewidth}}}]{
                \texttt{\tiny{[P1$\langle$GM]}}
                \texttt{Kazna: (7,6) nije prazno, već sadrži '─'.} \\
\\ 
\texttt{Zajedno ste prikupili 2 od 8 kazni.} \\
\texttt{Molim pokušajte ponovno!} \\
            }
        }
    }
    & \\ \\

    \theutterance \stepcounter{utterance}  
    & \multicolumn{4}{p{0.6\linewidth}}{
        \cellcolor[rgb]{0.9,0.9,0.9}{
            \makecell[{{p{\linewidth}}}]{
                \texttt{\tiny{[P1$\rangle$GM]}}
                \texttt{RECI: Oprosti na pogreškama! Čini se da se moji objekti već nalaze na ciljnim pozicijama koje si predložio (L na 1,1 i C na 1,2), pa mi ne treba pomicanje za njih. Za P, pretpostavljam da ćemo ga oboje pomicati prema nekoj zajedničkoj poziciji. Molim te, predloži gdje bi trebao biti P, a ja ću ga tamo pomaknuti. Također, provjeri možeš li ti pomicati svoje objekte na dogovorene pozicije.} \\
            }
        }
    }
    & & \\ \\

    \theutterance \stepcounter{utterance}  
    & & & \multicolumn{2}{p{0.3\linewidth}}{
        \cellcolor[rgb]{0.9,0.9,0.9}{
            \makecell[{{p{\linewidth}}}]{
                \texttt{\tiny{[GM$|$GM]}}
                \texttt{RECI: Oprosti na pogreškama! Čini se da se moji objekti već nalaze na ciljnim pozicijama koje si predložio (L na 1,1 i C na 1,2), pa mi ne treba pomicanje za njih. Za P, pretpostavljam da ćemo ga oboje pomicati prema nekoj zajedničkoj poziciji. Molim te, predloži gdje bi trebao biti P, a ja ću ga tamo pomaknuti. Također, provjeri možeš li ti pomicati svoje objekte na dogovorene pozicije.} \\
            }
        }
    }
    & & \\ \\

    \theutterance \stepcounter{utterance}  
    & & \multicolumn{4}{p{0.6\linewidth}}{
        \cellcolor[rgb]{0.9,0.9,0.9}{
            \makecell[{{p{\linewidth}}}]{
                \texttt{\tiny{[GM$\rangle$P2]}}
                \texttt{Vaša poruka je proslijeđena drugom igraču.} \\
\\ 
\\ 
\texttt{Trenutačno igrate rundu 2 od 12.} \\
\texttt{Zajedno ste prikupili 2 od 8 kazni.} \\
\texttt{Drugi igrač je poslao ovu poruku:} \\
\\ 
\\ 
\texttt{"Oprosti na pogreškama! Čini se da se moji objekti već nalaze na ciljnim pozicijama koje si predložio (L na 1,1 i C na 1,2), pa mi ne treba pomicanje za njih. Za P, pretpostavljam da ćemo ga oboje pomicati prema nekoj zajedničkoj poziciji. Molim te, predloži gdje bi trebao biti P, a ja ću ga tamo pomaknuti. Također, provjeri možeš li ti pomicati svoje objekte na dogovorene pozicije."} \\
\\ 
\\ 
\texttt{Koja je vaša sljedeća naredba?} \\
            }
        }
    }
    & \\ \\

    \theutterance \stepcounter{utterance}  
    & & & \multicolumn{4}{p{0.6\linewidth}}{
        \cellcolor[rgb]{0.9,0.9,0.9}{
            \makecell[{{p{\linewidth}}}]{
                \texttt{\tiny{[GM$\langle$P2]}}
                \texttt{RECI: Slazem se. Za P predlazem cilj (7,7). Ja cu sada pomaknuti L na (1,1), a P na (7,7). C mora ostati na (4,5), jer je (1,2) kod mene blokirana linijom. Molim te, pomakni svoj P na (7,7) ako vec nije tamo.} \\
            }
        }
    }
     \\ \\

    \theutterance \stepcounter{utterance}  
    & & & \multicolumn{2}{p{0.3\linewidth}}{
        \cellcolor[rgb]{0.9,0.9,0.9}{
            \makecell[{{p{\linewidth}}}]{
                \texttt{\tiny{[GM$|$GM]}}
                \texttt{RECI: Slazem se. Za P predlazem cilj (7,7). Ja cu sada pomaknuti L na (1,1), a P na (7,7). C mora ostati na (4,5), jer je (1,2) kod mene blokirana linijom. Molim te, pomakni svoj P na (7,7) ako vec nije tamo.} \\
            }
        }
    }
    & & \\ \\

    \theutterance \stepcounter{utterance}  
    & & \multicolumn{4}{p{0.6\linewidth}}{
        \cellcolor[rgb]{0.9,0.9,0.9}{
            \makecell[{{p{\linewidth}}}]{
                \texttt{\tiny{[P1$\langle$GM]}}
                \texttt{Vaša poruka je proslijeđena drugom igraču.} \\
\\ 
\\ 
\texttt{Trenutačno igrate rundu 2 od 12.} \\
\texttt{Zajedno ste prikupili 2 od 8 kazni.} \\
\texttt{Drugi igrač je poslao ovu poruku:} \\
\\ 
\\ 
\texttt{"Slazem se. Za P predlazem cilj (7,7). Ja cu sada pomaknuti L na (1,1), a P na (7,7). C mora ostati na (4,5), jer je (1,2) kod mene blokirana linijom. Molim te, pomakni svoj P na (7,7) ako vec nije tamo."} \\
\\ 
\\ 
\texttt{Koja je vaša sljedeća naredba?} \\
            }
        }
    }
    & \\ \\

    \theutterance \stepcounter{utterance}  
    & \multicolumn{4}{p{0.6\linewidth}}{
        \cellcolor[rgb]{0.9,0.9,0.9}{
            \makecell[{{p{\linewidth}}}]{
                \texttt{\tiny{[P1$\rangle$GM]}}
                \texttt{RECI: Moj je P već na (7,7), dakle ne moram ga pomicati. No, sada se javlja problem: moj je C na (1,2), a tvoj je C zaglavljen na (4,5) i ne možeš ga pomaknuti do (1,2). Znači, moramo pronaći novu ciljnu poziciju za C koja je prazna na obje ploče i do koje oboje možemo pomaknuti svoje C. Molim te, provjeri na svojoj ploči je li pozicija (4,4) prazna i možeš li pomaknuti svoj C tamo. Ako da, ja ću pokušati pomaknuti svoj C na (4,4).} \\
            }
        }
    }
    & & \\ \\

    \theutterance \stepcounter{utterance}  
    & & & \multicolumn{2}{p{0.3\linewidth}}{
        \cellcolor[rgb]{0.9,0.9,0.9}{
            \makecell[{{p{\linewidth}}}]{
                \texttt{\tiny{[GM$|$GM]}}
                \texttt{RECI: Moj je P već na (7,7), dakle ne moram ga pomicati. No, sada se javlja problem: moj je C na (1,2), a tvoj je C zaglavljen na (4,5) i ne možeš ga pomaknuti do (1,2). Znači, moramo pronaći novu ciljnu poziciju za C koja je prazna na obje ploče i do koje oboje možemo pomaknuti svoje C. Molim te, provjeri na svojoj ploči je li pozicija (4,4) prazna i možeš li pomaknuti svoj C tamo. Ako da, ja ću pokušati pomaknuti svoj C na (4,4).} \\
            }
        }
    }
    & & \\ \\

    \theutterance \stepcounter{utterance}  
    & & \multicolumn{4}{p{0.6\linewidth}}{
        \cellcolor[rgb]{0.9,0.9,0.9}{
            \makecell[{{p{\linewidth}}}]{
                \texttt{\tiny{[GM$\rangle$P2]}}
                \texttt{Vaša poruka je proslijeđena drugom igraču.} \\
\\ 
\\ 
\texttt{Trenutačno igrate rundu 3 od 12.} \\
\texttt{Zajedno ste prikupili 2 od 8 kazni.} \\
\texttt{Drugi igrač je poslao ovu poruku:} \\
\\ 
\\ 
\texttt{"Moj je P već na (7,7), dakle ne moram ga pomicati. No, sada se javlja problem: moj je C na (1,2), a tvoj je C zaglavljen na (4,5) i ne možeš ga pomaknuti do (1,2). Znači, moramo pronaći novu ciljnu poziciju za C koja je prazna na obje ploče i do koje oboje možemo pomaknuti svoje C. Molim te, provjeri na svojoj ploči je li pozicija (4,4) prazna i možeš li pomaknuti svoj C tamo. Ako da, ja ću pokušati pomaknuti svoj C na (4,4)."} \\
\\ 
\\ 
\texttt{Koja je vaša sljedeća naredba?} \\
            }
        }
    }
    & \\ \\

    \theutterance \stepcounter{utterance}  
    & & & \multicolumn{4}{p{0.6\linewidth}}{
        \cellcolor[rgb]{0.9,0.9,0.9}{
            \makecell[{{p{\linewidth}}}]{
                \texttt{\tiny{[GM$\langle$P2]}}
                \texttt{POMAKNI: L, (1, 1)} \\
            }
        }
    }
     \\ \\

    \theutterance \stepcounter{utterance}  
    & & & \multicolumn{2}{p{0.3\linewidth}}{
        \cellcolor[rgb]{0.9,0.9,0.9}{
            \makecell[{{p{\linewidth}}}]{
                \texttt{\tiny{[GM$|$GM]}}
                \texttt{POMAKNI: L, (1, 1)} \\
            }
        }
    }
    & & \\ \\

    \theutterance \stepcounter{utterance}  
    & & & \multicolumn{2}{p{0.3\linewidth}}{
        \cellcolor[rgb]{0.9,0.9,0.9}{
            \makecell[{{p{\linewidth}}}]{
                \texttt{\tiny{[GM$|$GM]}}
                \texttt{L moved to (1, 1)} \\
            }
        }
    }
    & & \\ \\

    \theutterance \stepcounter{utterance}  
    & & \multicolumn{4}{p{0.6\linewidth}}{
        \cellcolor[rgb]{0.9,0.9,0.9}{
            \makecell[{{p{\linewidth}}}]{
                \texttt{\tiny{[P1$\langle$GM]}}
                \texttt{Vaša poruka je proslijeđena drugom igraču.} \\
\\ 
\\ 
\texttt{Trenutačno igrate rundu 3 od 12.} \\
\texttt{Zajedno ste prikupili 2 od 8 kazni.} \\
\texttt{Drugi igrač je pomaknuo objekt na svojoj mreži.} \\
\\ 
\\ 
\texttt{Koja je vaša sljedeća naredba?} \\
            }
        }
    }
    & \\ \\

    \theutterance \stepcounter{utterance}  
    & \multicolumn{4}{p{0.6\linewidth}}{
        \cellcolor[rgb]{0.9,0.9,0.9}{
            \makecell[{{p{\linewidth}}}]{
                \texttt{\tiny{[P1$\rangle$GM]}}
                \texttt{POMAKNI: C, (4, 4)} \\
            }
        }
    }
    & & \\ \\

    \theutterance \stepcounter{utterance}  
    & & & \multicolumn{2}{p{0.3\linewidth}}{
        \cellcolor[rgb]{0.9,0.9,0.9}{
            \makecell[{{p{\linewidth}}}]{
                \texttt{\tiny{[GM$|$GM]}}
                \texttt{POMAKNI: C, (4, 4)} \\
            }
        }
    }
    & & \\ \\

    \theutterance \stepcounter{utterance}  
    & & & \multicolumn{2}{p{0.3\linewidth}}{
        \cellcolor[rgb]{0.9,0.9,0.9}{
            \makecell[{{p{\linewidth}}}]{
                \texttt{\tiny{[GM$|$GM]}}
                \texttt{Kazna: (4,4) nije prazno, već sadrži '│'.} \\
\\ 
\texttt{Zajedno ste prikupili 3 od 8 kazni.} \\
\texttt{Molim pokušajte ponovno!} \\
            }
        }
    }
    & & \\ \\

    \theutterance \stepcounter{utterance}  
    & & \multicolumn{4}{p{0.6\linewidth}}{
        \cellcolor[rgb]{0.9,0.9,0.9}{
            \makecell[{{p{\linewidth}}}]{
                \texttt{\tiny{[P1$\langle$GM]}}
                \texttt{Kazna: (4,4) nije prazno, već sadrži '│'.} \\
\\ 
\texttt{Zajedno ste prikupili 3 od 8 kazni.} \\
\texttt{Molim pokušajte ponovno!} \\
            }
        }
    }
    & \\ \\

    \theutterance \stepcounter{utterance}  
    & \multicolumn{4}{p{0.6\linewidth}}{
        \cellcolor[rgb]{0.9,0.9,0.9}{
            \makecell[{{p{\linewidth}}}]{
                \texttt{\tiny{[P1$\rangle$GM]}}
                \texttt{RECI: Pokušao sam pomaknuti C na (4,4), ali ta ćelija kod mene sadrži liniju '│', dakle nije prazna. Čini se da ploče nisu identične kad su u pitanju linije. Predlažem da si prekrižimo sve prazne ćelije koje su sigurno dostupne na obje ploče. Molim te, provjeri na svojoj ploči jesu li prazne i dostupne slobodne ćelije npr. na ovim pozicijama: (2,4), (3,4), (5,4), (6,4), (7,4), (2,5), (3,5), (4,6), (5,6), (6,6), (7,6), (2,7), (3,7), (4,7), (5,7), (6,7). Javi mi koje od njih mogu primiti C.} \\
            }
        }
    }
    & & \\ \\

    \theutterance \stepcounter{utterance}  
    & & & \multicolumn{2}{p{0.3\linewidth}}{
        \cellcolor[rgb]{0.9,0.9,0.9}{
            \makecell[{{p{\linewidth}}}]{
                \texttt{\tiny{[GM$|$GM]}}
                \texttt{RECI: Pokušao sam pomaknuti C na (4,4), ali ta ćelija kod mene sadrži liniju '│', dakle nije prazna. Čini se da ploče nisu identične kad su u pitanju linije. Predlažem da si prekrižimo sve prazne ćelije koje su sigurno dostupne na obje ploče. Molim te, provjeri na svojoj ploči jesu li prazne i dostupne slobodne ćelije npr. na ovim pozicijama: (2,4), (3,4), (5,4), (6,4), (7,4), (2,5), (3,5), (4,6), (5,6), (6,6), (7,6), (2,7), (3,7), (4,7), (5,7), (6,7). Javi mi koje od njih mogu primiti C.} \\
            }
        }
    }
    & & \\ \\

    \theutterance \stepcounter{utterance}  
    & & \multicolumn{4}{p{0.6\linewidth}}{
        \cellcolor[rgb]{0.9,0.9,0.9}{
            \makecell[{{p{\linewidth}}}]{
                \texttt{\tiny{[GM$\rangle$P2]}}
                \texttt{Uspješno premješteno 'L' na (1,1). Vaša ažurirana mreža izgleda ovako:} \\
\texttt{```} \\
\texttt{ 1234567} \\
\texttt{╔═╤═╤═══╗} \\
\texttt{║L│◌│◌◌◌║ 1} \\
\texttt{║◌├─┤◌◌◌║ 2} \\
\texttt{╟─┤P├─┬─╢ 3} \\
\texttt{║◌│◌│◌│◌║ 4} \\
\texttt{╟─┼─┘C│◌║ 5} \\
\texttt{║◌│◌◌┌┴─╢ 6} \\
\texttt{║◌│◌◌│◌◌║ 7} \\
\texttt{╚═╧══╧══╝} \\
\texttt{```} \\
\texttt{Trenutačno igrate rundu 4 od 12.} \\
\texttt{Zajedno ste prikupili 3 od 8 kazni.} \\
\texttt{Drugi igrač je poslao ovu poruku:} \\
\\ 
\\ 
\texttt{"Pokušao sam pomaknuti C na (4,4), ali ta ćelija kod mene sadrži liniju '│', dakle nije prazna. Čini se da ploče nisu identične kad su u pitanju linije. Predlažem da si prekrižimo sve prazne ćelije koje su sigurno dostupne na obje ploče. Molim te, provjeri na svojoj ploči jesu li prazne i dostupne slobodne ćelije npr. na ovim pozicijama: (2,4), (3,4), (5,4), (6,4), (7,4), (2,5), (3,5), (4,6), (5,6), (6,6), (7,6), (2,7), (3,7), (4,7), (5,7), (6,7). Javi mi koje od njih mogu primiti C."} \\
\\ 
\\ 
\texttt{Koja je vaša sljedeća naredba?} \\
            }
        }
    }
    & \\ \\

    \theutterance \stepcounter{utterance}  
    & & & \multicolumn{4}{p{0.6\linewidth}}{
        \cellcolor[rgb]{0.9,0.9,0.9}{
            \makecell[{{p{\linewidth}}}]{
                \texttt{\tiny{[GM$\langle$P2]}}
                \texttt{RECI: Pozicija (4,4) je prazna kod mene. Provjeri mozes li svoj C pomaknuti na (2,4), (3,4) ili (4,6) {-} koje od tih pozicija je prazno i dostupno kod tebe?} \\
            }
        }
    }
     \\ \\

    \theutterance \stepcounter{utterance}  
    & & & \multicolumn{2}{p{0.3\linewidth}}{
        \cellcolor[rgb]{0.9,0.9,0.9}{
            \makecell[{{p{\linewidth}}}]{
                \texttt{\tiny{[GM$|$GM]}}
                \texttt{RECI: Pozicija (4,4) je prazna kod mene. Provjeri mozes li svoj C pomaknuti na (2,4), (3,4) ili (4,6) {-} koje od tih pozicija je prazno i dostupno kod tebe?} \\
            }
        }
    }
    & & \\ \\

    \theutterance \stepcounter{utterance}  
    & & \multicolumn{4}{p{0.6\linewidth}}{
        \cellcolor[rgb]{0.9,0.9,0.9}{
            \makecell[{{p{\linewidth}}}]{
                \texttt{\tiny{[P1$\langle$GM]}}
                \texttt{Vaša poruka je proslijeđena drugom igraču.} \\
\\ 
\\ 
\texttt{Trenutačno igrate rundu 4 od 12.} \\
\texttt{Zajedno ste prikupili 3 od 8 kazni.} \\
\texttt{Drugi igrač je poslao ovu poruku:} \\
\\ 
\\ 
\texttt{"Pozicija (4,4) je prazna kod mene. Provjeri mozes li svoj C pomaknuti na (2,4), (3,4) ili (4,6) {-} koje od tih pozicija je prazno i dostupno kod tebe?"} \\
\\ 
\\ 
\texttt{Koja je vaša sljedeća naredba?} \\
            }
        }
    }
    & \\ \\

    \theutterance \stepcounter{utterance}  
    & \multicolumn{4}{p{0.6\linewidth}}{
        \cellcolor[rgb]{0.9,0.9,0.9}{
            \makecell[{{p{\linewidth}}}]{
                \texttt{\tiny{[P1$\rangle$GM]}}
                \texttt{RECI: Kod mene su prazne i dostupne pozicije (2,4), (3,4) i (4,6) za C. Predlažem da C pomaknemo na (4,6). Možeš li ti pomaknuti svoj C na (4,6)? Ako da, ja ću u sljedećem potezu pomaknuti svoj C tamo.} \\
            }
        }
    }
    & & \\ \\

    \theutterance \stepcounter{utterance}  
    & & & \multicolumn{2}{p{0.3\linewidth}}{
        \cellcolor[rgb]{0.9,0.9,0.9}{
            \makecell[{{p{\linewidth}}}]{
                \texttt{\tiny{[GM$|$GM]}}
                \texttt{RECI: Kod mene su prazne i dostupne pozicije (2,4), (3,4) i (4,6) za C. Predlažem da C pomaknemo na (4,6). Možeš li ti pomaknuti svoj C na (4,6)? Ako da, ja ću u sljedećem potezu pomaknuti svoj C tamo.} \\
            }
        }
    }
    & & \\ \\

    \theutterance \stepcounter{utterance}  
    & & \multicolumn{4}{p{0.6\linewidth}}{
        \cellcolor[rgb]{0.9,0.9,0.9}{
            \makecell[{{p{\linewidth}}}]{
                \texttt{\tiny{[GM$\rangle$P2]}}
                \texttt{Vaša poruka je proslijeđena drugom igraču.} \\
\\ 
\\ 
\texttt{Trenutačno igrate rundu 5 od 12.} \\
\texttt{Zajedno ste prikupili 3 od 8 kazni.} \\
\texttt{Drugi igrač je poslao ovu poruku:} \\
\\ 
\\ 
\texttt{"Kod mene su prazne i dostupne pozicije (2,4), (3,4) i (4,6) za C. Predlažem da C pomaknemo na (4,6). Možeš li ti pomaknuti svoj C na (4,6)? Ako da, ja ću u sljedećem potezu pomaknuti svoj C tamo."} \\
\\ 
\\ 
\texttt{Koja je vaša sljedeća naredba?} \\
            }
        }
    }
    & \\ \\

    \theutterance \stepcounter{utterance}  
    & & & \multicolumn{4}{p{0.6\linewidth}}{
        \cellcolor[rgb]{0.9,0.9,0.9}{
            \makecell[{{p{\linewidth}}}]{
                \texttt{\tiny{[GM$\langle$P2]}}
                \texttt{POMAKNI: P, (7, 7)} \\
            }
        }
    }
     \\ \\

    \theutterance \stepcounter{utterance}  
    & & & \multicolumn{2}{p{0.3\linewidth}}{
        \cellcolor[rgb]{0.9,0.9,0.9}{
            \makecell[{{p{\linewidth}}}]{
                \texttt{\tiny{[GM$|$GM]}}
                \texttt{POMAKNI: P, (7, 7)} \\
            }
        }
    }
    & & \\ \\

    \theutterance \stepcounter{utterance}  
    & & & \multicolumn{2}{p{0.3\linewidth}}{
        \cellcolor[rgb]{0.9,0.9,0.9}{
            \makecell[{{p{\linewidth}}}]{
                \texttt{\tiny{[GM$|$GM]}}
                \texttt{P moved to (7, 7)} \\
            }
        }
    }
    & & \\ \\

    \theutterance \stepcounter{utterance}  
    & & \multicolumn{4}{p{0.6\linewidth}}{
        \cellcolor[rgb]{0.9,0.9,0.9}{
            \makecell[{{p{\linewidth}}}]{
                \texttt{\tiny{[P1$\langle$GM]}}
                \texttt{Vaša poruka je proslijeđena drugom igraču.} \\
\\ 
\\ 
\texttt{Trenutačno igrate rundu 5 od 12.} \\
\texttt{Zajedno ste prikupili 3 od 8 kazni.} \\
\texttt{Drugi igrač je pomaknuo objekt na svojoj mreži.} \\
\\ 
\\ 
\texttt{Koja je vaša sljedeća naredba?} \\
            }
        }
    }
    & \\ \\

    \theutterance \stepcounter{utterance}  
    & \multicolumn{4}{p{0.6\linewidth}}{
        \cellcolor[rgb]{0.9,0.9,0.9}{
            \makecell[{{p{\linewidth}}}]{
                \texttt{\tiny{[P1$\rangle$GM]}}
                \texttt{POMAKNI: C, (4, 6)} \\
            }
        }
    }
    & & \\ \\

    \theutterance \stepcounter{utterance}  
    & & & \multicolumn{2}{p{0.3\linewidth}}{
        \cellcolor[rgb]{0.9,0.9,0.9}{
            \makecell[{{p{\linewidth}}}]{
                \texttt{\tiny{[GM$|$GM]}}
                \texttt{POMAKNI: C, (4, 6)} \\
            }
        }
    }
    & & \\ \\

    \theutterance \stepcounter{utterance}  
    & & & \multicolumn{2}{p{0.3\linewidth}}{
        \cellcolor[rgb]{0.9,0.9,0.9}{
            \makecell[{{p{\linewidth}}}]{
                \texttt{\tiny{[GM$|$GM]}}
                \texttt{C moved to (4, 6)} \\
            }
        }
    }
    & & \\ \\

    \theutterance \stepcounter{utterance}  
    & & \multicolumn{4}{p{0.6\linewidth}}{
        \cellcolor[rgb]{0.9,0.9,0.9}{
            \makecell[{{p{\linewidth}}}]{
                \texttt{\tiny{[GM$\rangle$P2]}}
                \texttt{Uspješno premješteno 'P' na (7,7). Vaša ažurirana mreža izgleda ovako:} \\
\texttt{```} \\
\texttt{ 1234567} \\
\texttt{╔═╤═╤═══╗} \\
\texttt{║L│◌│◌◌◌║ 1} \\
\texttt{║◌├─┤◌◌◌║ 2} \\
\texttt{╟─┤◌├─┬─╢ 3} \\
\texttt{║◌│◌│◌│◌║ 4} \\
\texttt{╟─┼─┘C│◌║ 5} \\
\texttt{║◌│◌◌┌┴─╢ 6} \\
\texttt{║◌│◌◌│◌P║ 7} \\
\texttt{╚═╧══╧══╝} \\
\texttt{```} \\
\texttt{Trenutačno igrate rundu 6 od 12.} \\
\texttt{Zajedno ste prikupili 3 od 8 kazni.} \\
\texttt{Drugi igrač je pomaknuo objekt na svojoj mreži.} \\
\\ 
\\ 
\texttt{Koja je vaša sljedeća naredba?} \\
            }
        }
    }
    & \\ \\

    \theutterance \stepcounter{utterance}  
    & & & \multicolumn{4}{p{0.6\linewidth}}{
        \cellcolor[rgb]{0.9,0.9,0.9}{
            \makecell[{{p{\linewidth}}}]{
                \texttt{\tiny{[GM$\langle$P2]}}
                \texttt{RECI: Pomaknuo sam P na (7,7). Sljedeći potez: moj C ide na (4,6). Javi mi kad pomakneš svoj C na (4,6).} \\
            }
        }
    }
     \\ \\

    \theutterance \stepcounter{utterance}  
    & & & \multicolumn{2}{p{0.3\linewidth}}{
        \cellcolor[rgb]{0.9,0.9,0.9}{
            \makecell[{{p{\linewidth}}}]{
                \texttt{\tiny{[GM$|$GM]}}
                \texttt{RECI: Pomaknuo sam P na (7,7). Sljedeći potez: moj C ide na (4,6). Javi mi kad pomakneš svoj C na (4,6).} \\
            }
        }
    }
    & & \\ \\

    \theutterance \stepcounter{utterance}  
    & & \multicolumn{4}{p{0.6\linewidth}}{
        \cellcolor[rgb]{0.9,0.9,0.9}{
            \makecell[{{p{\linewidth}}}]{
                \texttt{\tiny{[P1$\langle$GM]}}
                \texttt{Uspješno premješteno 'C' na (4,6). Vaša ažurirana mreža izgleda ovako:} \\
\texttt{```} \\
\texttt{ 1234567} \\
\texttt{╔═╤═╤═══╗} \\
\texttt{║L│◌│◌◌◌║ 1} \\
\texttt{║◌├─┤◌◌◌║ 2} \\
\texttt{╟─┤◌├─┬─╢ 3} \\
\texttt{║◌│◌│◌│◌║ 4} \\
\texttt{╟─┼─┘◌│◌║ 5} \\
\texttt{║◌│◌C┌┴─╢ 6} \\
\texttt{║◌│◌◌│◌P║ 7} \\
\texttt{╚═╧══╧══╝} \\
\texttt{```} \\
\texttt{Trenutačno igrate rundu 6 od 12.} \\
\texttt{Zajedno ste prikupili 3 od 8 kazni.} \\
\texttt{Drugi igrač je poslao ovu poruku:} \\
\\ 
\\ 
\texttt{"Pomaknuo sam P na (7,7). Sljedeći potez: moj C ide na (4,6). Javi mi kad pomakneš svoj C na (4,6)."} \\
\\ 
\\ 
\texttt{Koja je vaša sljedeća naredba?} \\
            }
        }
    }
    & \\ \\

    \theutterance \stepcounter{utterance}  
    & \multicolumn{4}{p{0.6\linewidth}}{
        \cellcolor[rgb]{0.9,0.9,0.9}{
            \makecell[{{p{\linewidth}}}]{
                \texttt{\tiny{[P1$\rangle$GM]}}
                \texttt{RECI: Upravo sam pomaknuo svoj C na (4,6), kao što smo se dogovorili. Sada su svi moji objekti na sljedećim koordinatama: L na (1,1), C na (4,6) i P na (7,7). Molim te, potvrdi kad i ti pomakneš svoj C na (4,6) te javi jesu li svi tvoji objekti na istim pozicijama. Ako je sve usklađeno, možemo završiti igru!} \\
            }
        }
    }
    & & \\ \\

    \theutterance \stepcounter{utterance}  
    & & & \multicolumn{2}{p{0.3\linewidth}}{
        \cellcolor[rgb]{0.9,0.9,0.9}{
            \makecell[{{p{\linewidth}}}]{
                \texttt{\tiny{[GM$|$GM]}}
                \texttt{RECI: Upravo sam pomaknuo svoj C na (4,6), kao što smo se dogovorili. Sada su svi moji objekti na sljedećim koordinatama: L na (1,1), C na (4,6) i P na (7,7). Molim te, potvrdi kad i ti pomakneš svoj C na (4,6) te javi jesu li svi tvoji objekti na istim pozicijama. Ako je sve usklađeno, možemo završiti igru!} \\
            }
        }
    }
    & & \\ \\

    \theutterance \stepcounter{utterance}  
    & & \multicolumn{4}{p{0.6\linewidth}}{
        \cellcolor[rgb]{0.9,0.9,0.9}{
            \makecell[{{p{\linewidth}}}]{
                \texttt{\tiny{[GM$\rangle$P2]}}
                \texttt{Vaša poruka je proslijeđena drugom igraču.} \\
\\ 
\\ 
\texttt{Trenutačno igrate rundu 7 od 12.} \\
\texttt{Zajedno ste prikupili 3 od 8 kazni.} \\
\texttt{Drugi igrač je poslao ovu poruku:} \\
\\ 
\\ 
\texttt{"Upravo sam pomaknuo svoj C na (4,6), kao što smo se dogovorili. Sada su svi moji objekti na sljedećim koordinatama: L na (1,1), C na (4,6) i P na (7,7). Molim te, potvrdi kad i ti pomakneš svoj C na (4,6) te javi jesu li svi tvoji objekti na istim pozicijama. Ako je sve usklađeno, možemo završiti igru!"} \\
\\ 
\\ 
\texttt{Koja je vaša sljedeća naredba?} \\
            }
        }
    }
    & \\ \\

    \theutterance \stepcounter{utterance}  
    & & & \multicolumn{4}{p{0.6\linewidth}}{
        \cellcolor[rgb]{0.9,0.9,0.9}{
            \makecell[{{p{\linewidth}}}]{
                \texttt{\tiny{[GM$\langle$P2]}}
                \texttt{POMAKNI: C, (4, 6)} \\
            }
        }
    }
     \\ \\

    \theutterance \stepcounter{utterance}  
    & & & \multicolumn{2}{p{0.3\linewidth}}{
        \cellcolor[rgb]{0.9,0.9,0.9}{
            \makecell[{{p{\linewidth}}}]{
                \texttt{\tiny{[GM$|$GM]}}
                \texttt{POMAKNI: C, (4, 6)} \\
            }
        }
    }
    & & \\ \\

    \theutterance \stepcounter{utterance}  
    & & & \multicolumn{2}{p{0.3\linewidth}}{
        \cellcolor[rgb]{0.9,0.9,0.9}{
            \makecell[{{p{\linewidth}}}]{
                \texttt{\tiny{[GM$|$GM]}}
                \texttt{C moved to (4, 6)} \\
            }
        }
    }
    & & \\ \\

    \theutterance \stepcounter{utterance}  
    & & \multicolumn{4}{p{0.6\linewidth}}{
        \cellcolor[rgb]{0.9,0.9,0.9}{
            \makecell[{{p{\linewidth}}}]{
                \texttt{\tiny{[P1$\langle$GM]}}
                \texttt{Vaša poruka je proslijeđena drugom igraču.} \\
\\ 
\\ 
\texttt{Trenutačno igrate rundu 7 od 12.} \\
\texttt{Zajedno ste prikupili 3 od 8 kazni.} \\
\texttt{Drugi igrač je pomaknuo objekt na svojoj mreži.} \\
\\ 
\\ 
\texttt{Koja je vaša sljedeća naredba?} \\
            }
        }
    }
    & \\ \\

    \theutterance \stepcounter{utterance}  
    & \multicolumn{4}{p{0.6\linewidth}}{
        \cellcolor[rgb]{0.9,0.9,0.9}{
            \makecell[{{p{\linewidth}}}]{
                \texttt{\tiny{[P1$\rangle$GM]}}
                \texttt{RECI: gotovo?} \\
            }
        }
    }
    & & \\ \\

    \theutterance \stepcounter{utterance}  
    & & & \multicolumn{2}{p{0.3\linewidth}}{
        \cellcolor[rgb]{0.9,0.9,0.9}{
            \makecell[{{p{\linewidth}}}]{
                \texttt{\tiny{[GM$|$GM]}}
                \texttt{RECI: gotovo?} \\
            }
        }
    }
    & & \\ \\

    \theutterance \stepcounter{utterance}  
    & & \multicolumn{4}{p{0.6\linewidth}}{
        \cellcolor[rgb]{0.9,0.9,0.9}{
            \makecell[{{p{\linewidth}}}]{
                \texttt{\tiny{[GM$\rangle$P2]}}
                \texttt{Uspješno premješteno 'C' na (4,6). Vaša ažurirana mreža izgleda ovako:} \\
\texttt{```} \\
\texttt{ 1234567} \\
\texttt{╔═╤═╤═══╗} \\
\texttt{║L│◌│◌◌◌║ 1} \\
\texttt{║◌├─┤◌◌◌║ 2} \\
\texttt{╟─┤◌├─┬─╢ 3} \\
\texttt{║◌│◌│◌│◌║ 4} \\
\texttt{╟─┼─┘◌│◌║ 5} \\
\texttt{║◌│◌C┌┴─╢ 6} \\
\texttt{║◌│◌◌│◌P║ 7} \\
\texttt{╚═╧══╧══╝} \\
\texttt{```} \\
\texttt{Trenutačno igrate rundu 8 od 12.} \\
\texttt{Zajedno ste prikupili 3 od 8 kazni.} \\
\texttt{Drugi igrač je poslao ovu poruku:} \\
\\ 
\\ 
\texttt{"gotovo?"} \\
\\ 
\\ 
\texttt{Koja je vaša sljedeća naredba?} \\
            }
        }
    }
    & \\ \\

    \theutterance \stepcounter{utterance}  
    & & & \multicolumn{4}{p{0.6\linewidth}}{
        \cellcolor[rgb]{0.9,0.9,0.9}{
            \makecell[{{p{\linewidth}}}]{
                \texttt{\tiny{[GM$\langle$P2]}}
                \texttt{RECI: gotovo!} \\
            }
        }
    }
     \\ \\

    \theutterance \stepcounter{utterance}  
    & & & \multicolumn{2}{p{0.3\linewidth}}{
        \cellcolor[rgb]{0.9,0.9,0.9}{
            \makecell[{{p{\linewidth}}}]{
                \texttt{\tiny{[GM$|$GM]}}
                \texttt{RECI: gotovo!} \\
            }
        }
    }
    & & \\ \\

    \theutterance \stepcounter{utterance}  
    & & & \multicolumn{2}{p{0.3\linewidth}}{
        \cellcolor[rgb]{0.9,0.9,0.9}{
            \makecell[{{p{\linewidth}}}]{
                \texttt{\tiny{[GM$|$GM]}}
                \texttt{true} \\
            }
        }
    }
    & & \\ \\

    \theutterance \stepcounter{utterance}  
    & & & \multicolumn{2}{p{0.3\linewidth}}{
        \cellcolor[rgb]{0.9,0.9,0.9}{
            \makecell[{{p{\linewidth}}}]{
                \texttt{\tiny{[GM$|$GM]}}
                \texttt{Initial states:} \\
\texttt{```} \\
\texttt{Player 1:} \\
\texttt{ 1234567} \\
\texttt{╔═╤═╤═══╗} \\
\texttt{║L│◌│◌◌◌║ 1} \\
\texttt{║C├─┤◌◌◌║ 2} \\
\texttt{╟─┤◌├─┬─╢ 3} \\
\texttt{║◌│◌│◌│◌║ 4} \\
\texttt{╟─┼─┘◌│◌║ 5} \\
\texttt{║◌│◌◌┌┴─╢ 6} \\
\texttt{║◌│◌◌│◌P║ 7} \\
\texttt{╚═╧══╧══╝} \\
\\ 
\\ 
\texttt{Player 2:} \\
\texttt{ 1234567} \\
\texttt{╔═╤═╤═══╗} \\
\texttt{║◌│◌│◌◌◌║ 1} \\
\texttt{║◌├─┤◌L◌║ 2} \\
\texttt{╟─┤P├─┬─╢ 3} \\
\texttt{║◌│◌│◌│◌║ 4} \\
\texttt{╟─┼─┘C│◌║ 5} \\
\texttt{║◌│◌◌┌┴─╢ 6} \\
\texttt{║◌│◌◌│◌◌║ 7} \\
\texttt{╚═╧══╧══╝} \\
\\ 
\texttt{```} \\
            }
        }
    }
    & & \\ \\

    \theutterance \stepcounter{utterance}  
    & & & \multicolumn{2}{p{0.3\linewidth}}{
        \cellcolor[rgb]{0.9,0.9,0.9}{
            \makecell[{{p{\linewidth}}}]{
                \texttt{\tiny{[GM$|$GM]}}
                \texttt{End states:} \\
\texttt{```} \\
\texttt{Player 1:} \\
\texttt{ 1234567} \\
\texttt{╔═╤═╤═══╗} \\
\texttt{║L│◌│◌◌◌║ 1} \\
\texttt{║◌├─┤◌◌◌║ 2} \\
\texttt{╟─┤◌├─┬─╢ 3} \\
\texttt{║◌│◌│◌│◌║ 4} \\
\texttt{╟─┼─┘◌│◌║ 5} \\
\texttt{║◌│◌C┌┴─╢ 6} \\
\texttt{║◌│◌◌│◌P║ 7} \\
\texttt{╚═╧══╧══╝} \\
\\ 
\\ 
\texttt{Player 2:} \\
\texttt{ 1234567} \\
\texttt{╔═╤═╤═══╗} \\
\texttt{║L│◌│◌◌◌║ 1} \\
\texttt{║◌├─┤◌◌◌║ 2} \\
\texttt{╟─┤◌├─┬─╢ 3} \\
\texttt{║◌│◌│◌│◌║ 4} \\
\texttt{╟─┼─┘◌│◌║ 5} \\
\texttt{║◌│◌C┌┴─╢ 6} \\
\texttt{║◌│◌◌│◌P║ 7} \\
\texttt{╚═╧══╧══╝} \\
\\ 
\texttt{```} \\
            }
        }
    }
    & & \\ \\

    \theutterance \stepcounter{utterance}  
    & & & \multicolumn{2}{p{0.3\linewidth}}{
        \cellcolor[rgb]{0.9,0.9,0.9}{
            \makecell[{{p{\linewidth}}}]{
                \texttt{\tiny{[GM$|$GM]}}
                \texttt{* success: True} \\
\texttt{* lose: False} \\
\texttt{* aborted: False} \\
\texttt{{-}{-}{-}{-}{-}{-}{-}} \\
\texttt{* Shifts: 2.00} \\
\texttt{* Max Shifts: 4.00} \\
\texttt{* Min Shifts: 2.00} \\
\texttt{* End Distance Sum: 0.00} \\
\texttt{* Init Distance Sum: 15.76} \\
\texttt{* Expected Distance Sum: 12.57} \\
\texttt{* Penalties: 3.00} \\
\texttt{* Max Penalties: 8.00} \\
\texttt{* Rounds: 8.00} \\
\texttt{* Max Rounds: 12.00} \\
\texttt{* Object Count: 3.00} \\
            }
        }
    }
    & & \\ \\

\end{supertabular}
}

\end{document}
